
\documentclass{article}
\usepackage{colortbl}
\usepackage{makecell}
\usepackage{multirow}
\usepackage{supertabular}

\begin{document}

\newcounter{utterance}

\centering \large Interaction Transcript for game `codenames', experiment `opponent\_none', episode 3 with mistral{-}large{-}3{-}azure.
\vspace{24pt}

{ \footnotesize  \setcounter{utterance}{1}
\setlength{\tabcolsep}{0pt}
\begin{supertabular}{c@{$\;$}|p{.15\linewidth}@{}p{.15\linewidth}p{.15\linewidth}p{.15\linewidth}p{.15\linewidth}p{.15\linewidth}}
    \# & $\;$A & \multicolumn{4}{c}{Game Master} & $\;\:$B\\
    \hline

    \theutterance \stepcounter{utterance}  
    & & & \multicolumn{2}{p{0.3\linewidth}}{
        \cellcolor[rgb]{0.9,0.9,0.9}{
            \makecell[{{p{\linewidth}}}]{
                \texttt{\tiny{[GM$|$GM]}}
                {'hidden': {'team': ['kenguru', 'saturn', 'nogavica', 'reševalno vozilo', 'vrstica', 'podkev', 'konica', 'učenec', 'čevelj'], 'innocent': ['tabla', 'trikotnik', 'stopalo', 'tableta', 'voda', 'stolp', 'oljka'], 'opponent': ['olimp', 'fant', 'jama', 'vrč', 'kocke', 'zrak', 'indija', 'london'], 'assassin': ['ninja']}, 'revealed': {'team': {'team': [], 'innocent': [], 'opponent': [], 'assassin': []}, 'opponent': {'team': [], 'innocent': [], 'opponent': [], 'assassin': []}}}
            }
        }
    }
    & & \\ \\

    \theutterance \stepcounter{utterance}  
    & & \multicolumn{4}{p{0.6\linewidth}}{
        \cellcolor[rgb]{0.9,0.9,0.9}{
            \makecell[{{p{\linewidth}}}]{
                \texttt{\tiny{[P1$\langle$GM]}}
                \texttt{Navedite eno enobesedno iztočnico, ki se nanaša na nekatere besede vaše ekipe s spodnjega seznama. Izberete lahko, na koliko in na katere besede se nanaša vaša iztočnica.} \\
\texttt{Cilj je, da vaš partner na podlagi vaših iztočnic ugane vse besede vaše ekipe, preden druga ekipa odkrije vse svoje besede. Če vaš partner ugane besedo "nasprotnika" ali "motečo" besedo, se njegovo ugibanje konča in na vrsto pride nasprotna ekipa. Če vaš partner ugane besedo "ubij", se igra konča in vaša ekipa takoj izgubi.} \\
\texttt{Več besed, na katere se nanaša vaša iztočnica, pomeni hitrejšo zmago v igri, vendar bo vaš partner morda težje uganil asociacije, zato pametno izberite iztočnico. Beseda iztočnica mora biti pomensko povezana s ciljnimi besedami, ne sme biti ena od besed na seznamih ali vsebovati njihovih delov.} \\
\\ 
\texttt{Vedno navedite enobesedno iztočnico in z vejico ločen seznam sorodnih ciljnih besed v naslednji obliki, odgovori naj bodo čim krajši, nikoli ne vključite drugega besedila, kot je zahtevano v tej obliki:} \\
\\ 
\texttt{NAMIG: <BESEDA>} \\
\texttt{CILJI: <CILJI>} \\
\\ 
\texttt{Vaše ekipne besede so: kenguru, saturn, nogavica, reševalno vozilo, vrstica, podkev, konica, učenec, čevelj.} \\
\texttt{Besede nasprotnikove ekipe so: olimp, fant, jama, vrč, kocke, zrak, indija, london.} \\
\texttt{Besede, ki odvračajo pozornost, so: tabla, trikotnik, stopalo, tableta, voda, stolp, oljka.} \\
\texttt{Besede za ubijanje so: ninja.} \\
            }
        }
    }
    & \\ \\

    \theutterance \stepcounter{utterance}  
    & \multicolumn{4}{p{0.6\linewidth}}{
        \cellcolor[rgb]{0.9,0.9,0.9}{
            \makecell[{{p{\linewidth}}}]{
                \texttt{\tiny{[P1$\rangle$GM]}}
                \texttt{NAMIG: žival} \\
\texttt{CILJI: kenguru, podkev} \\
\\ 
\texttt{NAMIG: obutev} \\
\texttt{CILJI: nogavica, čevelj} \\
\\ 
\texttt{NAMIG: konj} \\
\texttt{CILJI: podkev, konica} \\
\\ 
\texttt{NAMIG: šola} \\
\texttt{CILJI: vrstica, učenec} \\
\\ 
\texttt{NAMIG: planet} \\
\texttt{CILJI: saturn} \\
\\ 
\texttt{NAMIG: reševanje} \\
\texttt{CILJI: reševalno vozilo} \\
            }
        }
    }
    & & \\ \\

    \theutterance \stepcounter{utterance}  
    & & & \multicolumn{2}{p{0.3\linewidth}}{
        \cellcolor[rgb]{0.9,0.9,0.9}{
            \makecell[{{p{\linewidth}}}]{
                \texttt{\tiny{[GM$|$GM]}}
                {'player': 'Cluegiver', 'type': 'rambling error', 'utterance': 'NAMIG: žival\nCILJI: kenguru, podkev\n\nNAMIG: obutev\nCILJI: nogavica, čevelj\n\nNAMIG: konj\nCILJI: podkev, konica\n\nNAMIG: šola\nCILJI: vrstica, učenec\n\nNAMIG: planet\nCILJI: saturn\n\nNAMIG: reševanje\nCILJI: reševalno vozilo', 'message': 'Your answer contained more than two lines, please only give one clue and your targets on two separate lines.'}
            }
        }
    }
    & & \\ \\

\end{supertabular}
}

\end{document}
